\section{矩阵和线性映射}

记号约定:

\begin{itemize}
	\item $A$是幺环,$E,F,G$是左(相对的,右)$A$-模.
	\item $\mathcal{B}\coloneq\left(e_{i}\right)_{i\in I},\mathcal{C}\coloneq\left(c_{j}\right)_{j\in J},\mathcal{D}\coloneq\left(d_{k}\right)_{k\in K}$分别是$E,F,G$的基.
	\item $\mathcal{B}^{\ast}\coloneq\left(e_{i}^{\ast}\right)_{i\in I},\mathcal{C}^{\ast}\coloneq\left(c_{j}^{\ast}\right)_{j\in J}$分别是$\mathcal{B},\mathcal{C}$的对偶基.
\end{itemize}

\begin{definition}{向量的坐标矩阵}{}
	设$x\in E$.定义列矩阵$\left(x_{i0}\right)_{\left(i,0\right)\in I\times\left\{0\right\}}$如下:
	\begin{enumerate}
		\item 当$E$是左模时,$x=\sum\limits_{i\in I}x_{i0}e_{i}$.
		\item 当$E$是右模时,$x=\sum\limits_{i\in I}e_{i0}x_{i}$.
	\end{enumerate}
	我们称$\left(x_{i0}\right)_{\left(i,0\right)\in I\times\left\{0\right\}}$是$x$在$\mathcal{B}$下的矩阵.
\end{definition}

\begin{definition}{坐标映射}{}
	定义$\mathcal{B}$诱导的坐标映射$\left[\cdot\right]_{\mathcal{B}}:E\to R^{I\times\left\{0\right\}}$如下:对任一$x\in E$,$\left[x\right]_{\mathcal{B}}$是$x$在$\mathcal{B}$下的坐标矩阵.
\end{definition}

\begin{definition}{线性映射的表示矩阵}{}
	设$f\in\Hom_{A}\left(E,F\right)$,$A$上的矩阵$\left(a_{ji}\right)_{j\in J,i\in I}$满足$\forall i\in I\forall j\in J\left(a_{ji}=f_{j}^{\ast}\left(f\left(e_{i}\right)\right)\right)$.称$\left(a_{ji}\right)_{j\in J,i\in I}$是$f$在$\mathcal{B}$和$\mathcal{C}$下的表示矩阵.具体而言,记$\langle-,-\rangle_{F}$是$F$上的典范配对,那么:
	\begin{enumerate}
		\item 当$E,F$都是左$A$-模时,$\forall i\in I\forall j\in J\left(a_{j,i}=\langle f\left(e_{i}\right),f_{j}^{\ast}\rangle_{F}\right)$
		\item 当$E,F$都是右$A$-模时,$\forall i\in I\forall j\in J\left(a_{j,i}=\langle f_{j}^{\ast},f\left(e_{i}\right)\rangle_{F}\right)$.
	\end{enumerate}
\end{definition}

\begin{definition}{表示矩阵的转换映射}{}
	定义$\mathcal{B}$和$\mathcal{C}$诱导的转换映射$\left[\cdot\right]_{\mathcal{B}}^{\mathcal{C}}:\Hom_{A}\left(E,F\right)\to A^{J\times I}$如下:对任一$f\in\Hom_{A}\left(E,F\right)$,$\left[f\right]_{\mathcal{B}}^{\mathcal{C}}$是$f$在$\mathcal{B}$和$\mathcal{C}$下的表示矩阵.
\end{definition}

\begin{remark}{表示矩阵的每一列都是基在线性映射下的像的坐标矩阵}{}
	对任意$f\in\Hom_{A}\left(E,F\right)$和$i\in I$,有$\col_{i}\left(\left[f\right]_{\mathcal{B}}^{\mathcal{C}}\right)=\left[f\left(e_{i}\right)\right]_{\mathcal{C}}$.
\end{remark}

\begin{proposition}{转换映射是模同态}{}
	$\left[\cdot\right]_{\mathcal{B}}^{\mathcal{C}}$是$Z\left(A\right)$-模同态,并且在集合$J$有限时是$Z\left(A\right)$-模同构.
\end{proposition}

\begin{proposition}{映射的复合的表示矩阵}{}
	设集合$I,J,K$都有限.
	\begin{enumerate}
		\item 若$E,F,G$都是左$A$-模,则对任意$f\in\Hom_{A}\left(E,F\right),g\in\Hom_{A}\left(F,G\right)$,有$\left(\left[g\circ f\right]_{\mathcal{B}}^{\mathcal{D}}\right)^{\top}=\left(\left[f\right]_{\mathcal{B}}^{\mathcal{C}}\right)^{\top}\ast\left(\left[g\right]_{\mathcal{C}}^{\mathcal{D}}\right)^{\top}$.
		\item 若$E,F,G$都是右$A$-模,则对任意$f\in\Hom_{A}\left(E,F\right),g\in\Hom_{A}\left(F,G\right)$,有$\left[g\circ f\right]_{\mathcal{B}}^{\mathcal{D}}=\left[g\right]_{\mathcal{C}}^{\mathcal{D}}\left[f\right]_{\mathcal{B}}^{\mathcal{C}}$.
	\end{enumerate}
\end{proposition}

\begin{corollary}{映射作用下的坐标变换}{}
	设集合$I,J$都有限.
	\begin{enumerate}
		\item 若$E,F,G$都是左$A$-模,则对任意$f\in\Hom_{A}\left(E,F\right),x\in E$,有$\left(\left[f\left(x\right)\right]_{\mathcal{C}}\right)^{\top}=\left(\left[x\right]_{\mathcal{C}}\right)^{\top}\ast\left(\left[f\right]_{\mathcal{B}}^{\mathcal{C}}\right)^{\top}$.
		\item 若$E,F,G$都是右$A$-模,则对任意$f\in\Hom_{A}\left(E,F\right),x\in E$,有$\left[f\left(x\right)\right]_{\mathcal{C}}=\left[f\right]_{\mathcal{B}}^{\mathcal{C}}\left[x\right]_{\mathcal{C}}$.
	\end{enumerate}
\end{corollary}

\begin{proposition}{线性映射的对偶的表示矩阵}{}
	对任一$f\in\Hom_{A}\left(E,F\right)$,如下图表交换.
	\begin{center}
		\begin{tikzcd}
			{\Hom_{A}\left(E,F\right)} &&&& {\Hom_{A}\left(F^{\vee},E^{\vee}\right)} \\
			\\
			\\
			{A^{J\times I}} &&&& {A^{I\times J}}
			\arrow["{\Hom_{A}\left(f,\id_{F}\right)}", from=1-1, to=1-5]
			\arrow["{\left[\cdot\right]_{\mathcal{B}}^{\mathcal{C}}}", from=1-1, to=4-1]
			\arrow["{\left[\cdot\right]_{\mathcal{C^{\ast}}}^{\mathcal{B}^{\ast}}}"', from=1-5, to=4-5]
			\arrow["{\text{矩阵转置}}", from=4-1, to=4-5]
		\end{tikzcd}
	\end{center}
\end{proposition}

\begin{proposition}{标量扩张下线性映射的表示矩阵}{}
	设$A,B$是交换环,$\sigma\in\Hom_{\mathsf{Ring}}\left(A,B\right)$,集合$I,J$有限.记$\rho^{\ast}\left(E\right)$和$\rho^{\ast}\left(F\right)$各自的基是$\mathcal{B}^{\prime}$和$\mathcal{C}^{\prime}$,则对任一$f\in\Hom_{A}\left(E,F\right)$,
	\begin{align*}
		\left[\rho^{\ast}\left(f\right)\right]_{\mathcal{B}^{\prime}}^{\mathcal{C}^{\prime}}=\left(\left[f\right]_{\mathcal{B}}^{\mathcal{C}}\right)^{\sigma}.
	\end{align*}
\end{proposition}

\begin{definition}{线性方程组的表示矩阵}{matrix-of-linear-system}
	设集合$I,J$有限,$\left(a_{j,i}\right)_{j\in J,I\times I}$是$A$的矩阵,$\left(x_{i}\right)_{i\in I},\left(y_{j}\right)_{j\in J}$是$A$的元素族,$E=A_{d}^{\left(I\right)},F=A_{d}^{\left(J\right)}$,$\mathcal{B}$和$\mathcal{C}$分别是$E$和$F$的典范基,$f\in\Hom_{A}\left(E,F\right)$且$\boldsymbol{A}=\left[f\right]_{\mathcal{B}}^{\mathcal{C}}$.我们称$\boldsymbol{A}$是线性方程组$\sum\limits_{i\in I}a_{j,i}x_{i}=y_{j}$($j\in J$)的系数矩阵.
	\begin{align*}
		\boldsymbol{A}\coloneq\left[a_{j,i}\right]_{\substack{j\in J\\i\in I}}=\left[f\right]_{\mathcal{C}\leftarrow\mathcal{B}}.
	\end{align*}
	我们称$\boldsymbol{A}$是右标量线性方程组$\sum\limits_{i\in I}a_{j,i}x_{i}=y_{j}$($j\in J$)的矩阵.
\end{definition}

\begin{remark}{线性方程组的矩阵形式}{}
	记$\boldsymbol{Y}\coloneq\left(y_{j}\right)_{j\in J}$是列矩阵,则线性方程组$\sum\limits_{i\in I}a_{j,i}x_{i}=y_{j}$($j\in J$)等价于线性方程$\sum\limits_{i\in I}x_{i}\col_{i}\left(\boldsymbol{A}\right)=\boldsymbol{Y}$.
\end{remark}

\begin{theorem}{线性方程组有解的另一判别形式}{}
	沿用定义(\ref{def:matrix-of-linear-system})的记号,并记$\boldsymbol{Y}\coloneq\left(y_{j}\right)_{j\in J}$是$A$中的列矩阵,则
	\begin{center}
		线性方程组$\sum\limits_{i\in I}a_{j,i}x_{i}=y_{j}$($j\in J$)有解$\iff$ $\boldsymbol{Y}$是$\left(\col_{i}\left(\boldsymbol{A}\right)\right)_{i\in I}$的线性组合.
	\end{center}
\end{theorem}